\documentclass[11pt,a4paper]{article}
\usepackage[margin=26mm]{geometry}
\usepackage{amsmath,amssymb,amsthm,hyperref}
\newtheorem{lemma}{Lemma}
\newtheorem{theorem}[lemma]{Theorem}
\newcommand{\Alt}{\operatorname{Alt}}
\newcommand{\Sym}{\operatorname{Sym}}
\title{Alternating groups generated by a cycle and two transpositions}
\author{Human Compression project}\date{}
\begin{document}\maketitle
We give complete elementary generation and chirality arguments for the
two families stated in Section~5 of Leemans--Liebeck~\cite{LL}. The source
attributes these particular permutations to Gareth Jones. The families,
and the existence conclusion, are not new results of this note. Our purpose
is to supply their complete proof through one elementary generation lemma,
so that no search over the alternating groups is required.

\section{A generation lemma}
For a finite set $U$, let $\Alt(U)$ act trivially outside $U$.
\begin{lemma}\label{merge}
Suppose $|U|\geq3$ and $T$ is a set of three points meeting $U$.
Then $\langle\Alt(U),\Alt(T)\rangle=\Alt(U\cup T)$.
\end{lemma}
\begin{proof}
First suppose $T$ has two points $a,b$ in $U$ and one new point $d$.
Conjugating $(a\ b\ d)$ by $\Alt(U)$ gives a $3$-cycle supported on
$\{i,j,d\}$ for every pair of distinct points $i,j\in U$.
When $|U|\geq4$, this follows because $\Alt(U)$ is transitive on ordered
pairs: choose any permutation taking one ordered pair to the other and,
if its sign is wrong, compose with a transposition of two other points.
When $|U|=3$, its cyclic action is transitive on the three unordered pairs;
inverses give either orientation of the resulting $3$-cycles.
Together with $\Alt(U)$, these are all $3$-cycles on $U\cup\{d\}$.
They generate its alternating group. Indeed an even permutation is a product
of an even number of transpositions; pairing them expresses it as a product
of $3$-cycles, since two intersecting transpositions have a $3$-cycle as
product and, for distinct points, $(a\ b)(c\ d)=(a\ c\ b)(a\ c\ d)$.

If $T\cap U=\{a\}$, write $T=\{a,d,e\}$ and choose distinct
$b,c\in U\setminus\{a\}$. Conjugating $(a\ b\ c)$ by a suitable
element of $\Alt(T)$ gives a $3$-cycle on $\{d,b,c\}$.
The first case therefore supplies $\Alt(U\cup\{d\})$.
The original set $T$ now meets that enlarged set in two points, so a
second application supplies $\Alt(U\cup T)$. The case $T\subseteq U$
is immediate. All generators are even, giving equality rather than merely
containment.
\end{proof}

\begin{lemma}\label{primitive}
A primitive permutation group on a finite set $\Omega$ containing a
$3$-cycle contains $\Alt(\Omega)$.
\end{lemma}
\begin{proof}
Take the supports of all conjugates of that $3$-cycle. Regard these triples
as the edges of a hypergraph on $\Omega$. Its connected components form a
partition preserved by the group, because conjugation permutes the triples.
There is a component of size at least three. Transitivity and primitivity
therefore force the hypergraph to be connected. Start with one triple and
successively adjoin a triple meeting the union already obtained until all
vertices are included. Such an ordering exists by connectedness.
Each conjugate generates the alternating group of its three-point support.
Lemma~\ref{merge} consequently shows that the subgroup generated by all
these conjugates contains $\Alt(\Omega)$.
\end{proof}
This is the familiar $3$-cycle case of the primitive-group generation
criterion. The argument above includes its proof; it is not a new
classification theorem.

\section{The two families}
We use the standard fact $\operatorname{Aut}(\Alt(n))\cong\Sym(n)$ for
$n\geq8$, where a permutation acts by conjugation. This automorphism
theorem is retained as an external group-theoretic premise; we do not claim
to prove it here.

\begin{theorem}\label{pairs}
For even $n\geq8$, put
\[
 x=(2\ 3\ \cdots\ n),\qquad t=(1\ 2)(3\ 4).
\]
For odd $n\geq9$, put
\[
 x=(1\ 2\ \cdots\ n),\qquad t=(1\ 2)(3\ 6).
\]
In either case $\langle x,t\rangle=\Alt(n)$, and no automorphism of
$\Alt(n)$ both inverts $x$ and fixes $t$.
\end{theorem}
\begin{proof}
Let $G=\langle x,t\rangle$. Both permutations are even, so $G\leq\Alt(n)$.
In the even case, $x$ is transitive on all points except $1$, and $t$ moves
$1$ into that orbit. Hence $G$ is transitive. In any $G$-invariant partition,
the part containing $1$ is fixed by $x$. It therefore contains either none
or all of $\{2,\ldots,n\}$. The partition is trivial, proving primitivity.
The conjugate of $t$ obtained by applying $x^3$ to its point labels is
$(1\ 5)(6\ 7)$. Its product with $t$ has a $3$-cycle on $\{1,2,5\}$
and the two disjoint transpositions $(3\ 4),(6\ 7)$.
Squaring this product leaves a single $3$-cycle. Lemma~\ref{primitive}
now gives $G=\Alt(n)$.

In the odd case, $x$ already makes $G$ transitive. In a nontrivial invariant
partition, every part would have the same size, a proper divisor of the
odd number $n$, and hence size at least three. If $t$ moved one part to
another, it would move every point of at least two parts, and thus at
least six points. This contradicts its support of size four. Thus $t$
would fix every part, forcing $1$ and $2$ into the same part. Applying
powers of $x$ puts every consecutive pair in a common part and forces
the whole set into one part, a contradiction. Thus $G$ is primitive.
Here the conjugate obtained by applying $x^2$ to the labels of $t$ is
$(3\ 4)(5\ 8)$. Its product with $t$ consists of a $3$-cycle on
$\{3,4,6\}$ and the disjoint transpositions $(1\ 2),(5\ 8)$.
Its square is a $3$-cycle, so again Lemma~\ref{primitive} gives $G=\Alt(n)$.

It remains to exclude the unwanted automorphism. By the retained
automorphism theorem it is induced by a permutation $h$.
In the even case, inversion of $x$ makes $h$ fix its unique fixed point $1$.
If $h$ also fixes $t$ under conjugation, it must fix $2$, since $(1\ 2)$
is the unique transposition of $t$ meeting $1$. Reversing the cycle $x$
while fixing $2$ sends $3$ to $n$ and $4$ to $n-1$.
It cannot preserve the other transposition $(3\ 4)$ for $n\geq8$.

In the odd case, identify the point labels with $\mathbb Z/n\mathbb Z$
by subtracting one. Then $x$ is translation by one. Any permutation
inverting $x$ has the form $z\mapsto c-z$: its value at zero determines
all its other values by the inversion equation. The two transpositions
of $t$ have supports $\{0,1\}$ and $\{2,5\}$.
A reflection preserves their cyclic differences up to sign. Since
$1\not\equiv\pm3\pmod n$ for $n\geq9$, it cannot interchange these two
supports. Preserving them separately would require simultaneously
$c\equiv1$ and $c\equiv7\pmod n$, which is impossible. This completes
both chirality arguments.
\end{proof}

The usual correspondence between chiral polyhedra and generating pairs
$(x,t)$, with $t$ an involution and no inverting automorphism fixing $t$,
therefore supplies the alternating-group existence branch of the source
theorem. That correspondence is a retained input, attributed to
Schulte--Weiss as in the source. The theorem above proves the actual group
conditions. It does not assert the other classical or exceptional Lie
families, sporadic cases, their computer searches, or the almost-simple
question in the concluding section. Those obligations remain separate.

\begin{thebibliography}{9}
\bibitem{LL} D.~Leemans and M.~W.~Liebeck,
\emph{Chiral polyhedra and finite simple groups}, fixed retained
arXiv:1603.07713v3 PDF, arXiv header 18 May 2016, internal date
6 November 2018, Section~5. The displayed generating pairs are attributed
there to Gareth Jones's February 2016 Queenstown lecture.
\bibitem{SW} E.~Schulte and A.~Ivi\'c Weiss, \emph{Chiral polytopes},
Applied Geometry and Discrete Mathematics, DIMACS Series 4 (1991),493--516.
\end{thebibliography}
\end{document}
